% 05-conclusoes.tex | 5 Conclusions
% RESTAURADO DO ORIGINAL em 06/11/2026: secao "10. Conclusion" do Paper Training - 15 Nov 2025, na
% mesma estrutura (respostas as quatro perguntas, contribuicao, desafios e limitacoes, implicacoes
% para a pratica, trabalhos futuros, fechamento), condensada sem retirar nenhum ponto.
% Da versao JATM: a tabela de limitacoes (pedido 6 do revisor do JATM), ajustada ao desenho do
% original (Nivel 3 por questionarios pareados), e tres direcoes de pesquisa.

\section{Conclusions}
\label{sec:conclusions}

This research reframes systematic training evaluation from a regulatory obligation into a strategic priority that improves operational, financial and sustainability performance in aircraft engine \MRO{} at the same time. The answers to its four research questions follow, with theoretical and practical insights.

\textbf{RQ1.} The \AS{} and CAR training requirements form an integrated system that calls for a data-driven approach to competency. \AS{} requires objective evidence of competence, with documented processes for defining, delivering and evaluating training effectiveness. CARs, through the \SMS{} pillars, enforce a proactive safety culture in which training must be continuously evaluated to mitigate operational risk. Meeting these mandates requires more than record-keeping: it requires a structured evaluation framework.

\textbf{RQ2.} The training management process is optimized by embedding Kirkpatrick's four levels in the training life cycle, which turns a linear sequence (plan, deliver, archive) into an evidence-informed cycle. In this cycle, reaction data refine content and delivery, \LMS{} test scores expose knowledge gaps, trainee and leader evaluations verify on-the-job application, and customer and quality metrics align training with strategic outcomes.

\textbf{RQ3.} The model is operationalized by using the data of each level to feed a plan--do--check--act cycle that runs in step with the \AS{} quality system and the \SMS{}. Low scores at the behavior level, for instance, trigger root-cause analysis and content revision, providing auditable evidence of corrective action for both \AS{} and \CAA{} requirements. Evaluation thus directly informs training improvement, satisfying compliance while driving performance.

\textbf{RQ4.} The impact of ineffective training is severe and twofold. Financially, it inflates the cost of quality through internal failure costs (rework, scrap, delays) and external failure costs (warranty claims, concessions). In sustainability terms, these failures are direct indicators of resource waste: materials are scrapped, energy is consumed in reprocessing and human capital is underused. The study linked training deficiencies to substantial avoidable costs and environmental waste, indicating that effective training is a lever for cleaner production \parencite{Chen2023}.

\textbf{Contribution to the literature.} The findings provide empirical evidence of the Kirkpatrick model's usefulness in a high-risk aviation context, within the frame of sustainable operations. While prior research has explored the model's theoretical benefits, this case offers an application that identifies training deficiencies and supports improvements in both safety and resource efficiency. The result is consistent with studies that link human capital development to waste reduction in high-reliability organizations \parencite{Sousa2025} and shows the model's adaptability as a tool for cleaner production in industrial settings.

\textbf{Challenges and limitations.} Isolating the effects of training from other variables, such as concurrent changes in workplace policies or supply chain factors, is difficult, a common limitation in training evaluation research. The complexity of maintenance processes and the resources needed for evaluation may limit generalization to smaller organizations. The reliance on self-reported data at some levels introduces potential bias, although the mixed-methods approach, with objective operational and quality metrics, mitigates it. Table~\ref{tab:limitations} details each limitation, the mitigation adopted and the corresponding research direction.

% Tabela 2 | Limitacoes (origem: artigo/tabelas/table2-limitations.tex, versao JATM, pedido 6 do revisor;
% elogiada no parecer de 11/08/2026). Ajustada em 06/11/2026 ao desenho do ORIGINAL: o Nivel 3 usou
% questionarios pareados (trainee e lider), nao observacao estruturada; "workforce held constant" e
% "scenario bounds" sairam porque nao constam do original. Citacoes do original nao localizadas em vermelho.
\begin{table}[H]
  \centering
  \caption{Limitations of the study, mitigation adopted and future research directions}
  \label{tab:limitations}
  \footnotesize
  \begin{tabularx}{\linewidth}{@{}>{\raggedright\arraybackslash}p{3.6cm} >{\raggedright\arraybackslash}X >{\raggedright\arraybackslash}X@{}}
    \toprule
    \textbf{Limitation} & \textbf{Mitigation in this study} & \textbf{Future research direction} \\
    \midrule
    Single case (one \MRO{} facility) &
    Triangulation across four data sources (questionnaires, learning records, quality records, customer surveys) &
    Multi-site replication across regions and regulatory regimes (\FAA{}, \EASA{}, \ANAC{}) \\
    \addlinespace
    Training effects not isolated from other variables, such as concurrent policy or supply chain changes &
    Mixed methods; results read as associations &
    Quasi-experimental or difference-in-differences designs to separate contemporaneous confounders \parencite{FernandezDelgado2024,Patel2023} \\
    \addlinespace
    Self-reported behavior data &
    Paired trainee and leader questionnaires, complemented by objective quality and customer metrics &
    Digital or sensor-based measurement for continuous Level~3 evidence \parencite{Garcia2024} \\
    \addlinespace
    Generalizability to smaller \MRO{} facilities &
    Roll-out steps designed for integration into existing \LMS{} and \ERP{} infrastructure &
    Simplified framework for low-resource settings (lighter Level~3 protocol, sampled Level~4 metrics) \\
    \addlinespace
    Small, purposive samples (10 sessions, 21 trainees, 15 leader--trainee pairs, four customers) &
    Triangulation across levels and sources; results reported descriptively &
    Larger cohorts with sampling criteria and response rates reported \\
    \addlinespace
    No pre-intervention measure at Levels 1, 2 and 4 &
    Before-and-after comparison limited to training failures and regulatory non-conformances &
    Baseline measurement at every level before the next training cycle \\
    \addlinespace
    Sustainability outcomes described qualitatively, not yet in physical units &
    Cost-of-quality reasoning links avoided failures to material, labor and energy &
    Life-cycle assessment and Scope~3 accounting to turn cost of quality into environmental indicators \\
    \addlinespace
    Aggregated data only (no individual records released) &
    Anonymization and aggregation disclosed in the data availability statement and figure notes &
    Multi-organization consortium for shared anonymized data sets \\
    \addlinespace
    Insider access: one author is affiliated with the host organization and the team that implemented the framework also evaluated it, which gives privileged data but may introduce framing bias &
    External metrics verifiable outside the organization (authority audit outcomes, customer surveys); declaration in the conflict-of-interest statement &
    Independent replication in unaffiliated \MRO{} facilities \\
    \bottomrule
  \end{tabularx}
  \notafonte{MRO: maintenance, repair and overhaul; FAA: Federal Aviation Administration; EASA: European Union Aviation Safety Agency; ANAC: Brazilian National Civil Aviation Agency; LMS: learning management system; ERP: enterprise resource planning. Source: the authors.}
\end{table}


\textbf{Implications for practice.} Aviation maintenance organizations that adopt this integrated framework can systematically evaluate training programs, identify weaknesses and implement targeted improvements that enhance safety, compliance and resource efficiency at the same time. The structured approach supports continuous improvement and lets organizations adapt training to evolving standards and technologies. Regulators and industry associations may promote such frameworks to standardize evaluation practice and improve safety and environmental performance across the sector.

\textbf{Future research directions.} Future research should examine whether the improvements persist over multiple operational cycles, particularly in reducing material and energy waste. It should also investigate emerging technologies, such as artificial intelligence to predict the skill gaps that lead to waste or virtual reality to simulate resource-efficient procedures. Comparisons across aviation maintenance organizations and other high-risk industries would clarify the model's adaptability and scalability. Comparing regulatory regimes (\FAA{}, \EASA{} and \ANAC{}), measuring Level~3 continuously with digital tools to reduce reliance on self-report, and applying formal causal-attribution designs to longitudinal data would further strengthen the evidence.

By answering its research questions, the study shows that integrating Kirkpatrick's model within the aviation regulatory framework can turn training from a cost center into a value driver for safety, competitiveness and ecological stewardship. It gives \MRO{} organizations an actionable blueprint for compliance, safety and corporate sustainability goals, with the potential to reduce waste systematically and to make better use of human and material resources.
